\BourbakiExercises{Exercises}

\begin{enumerate}
\def\labelenumi{\arabic{enumi})}
\tightlist
\item
  Let \(\mathrm{K}\) be a commutative field, \(\mathrm{A}\) a central simple K-algebra, \(\mathrm{L}\) a semisimple commutative subalgebra of \(\mathrm{A}\) , and \(\mathrm{L}'\) its commutant in \(\mathrm{A}\) .
\end{enumerate}

\begin{enumerate}
\def\labelenumi{\alph{enumi})}
\item
  Suppose that A is the endomorphism algebra of a finite-dimensional K-vector space V. Prove that we have \([L:K][L':K] \geqslant [A:K]\) (cf.~VIII, p.~78, Proposition 1); equality holds if and only if the L-module V is free. Deduce examples where the inequality is strict.
\item
  Prove the inequality \([L:K][L':K] \geqslant [A:K]\) in general.
\end{enumerate}

\begin{enumerate}
\def\labelenumi{\arabic{enumi})}
\setcounter{enumi}{1}
\tightlist
\item
  Let K be a commutative field, A the K-algebra \(\mathbf{M}_{5}(\mathrm{K})\) , B the subalgebra of matrices of the form
\end{enumerate}

\[
\left( \begin{array}{c c c} x & 0 & 0 \\ 0 & Y & 0 \\ 0 & 0 & Y \end{array} \right)  ,
\]

and C the subalgebra of matrices of the form

\[
\left( \begin{array}{c c c c} x & 0 & 0 & 0 \\ 0 & x & 0 & 0 \\ 0 & 0 & x & 0 \\ 0 & 0 & 0 & Y \end{array} \right)  ,
\]

where x runs through K and Y through the algebra \(\mathbf{M}_{2}(\mathrm{K})\) . Prove that B and C are semisimple algebras containing the center Z of A. Then, prove that there exists an isomorphism of K-algebras \(f: B \to C\) that leaves invariant the elements of Z but that there does not exist any automorphism of A extending f (observe that such an automorphism must send a simple component of B to a simple component of C).

¶ 3) a) Let K be a commutative field, A a K-algebra, and B a central simple subalgebra of A of finite rank. Let B' be the commutant of B in A. Prove that the canonical homomorphism from B ⊗K B' to A is bijective. (Note that B and B' are linearly disjoint over K (VIII, p.~225, Exercise 5, d)). Prove, furthermore, that A is the direct sum of a family (M \(_{t}\) ) \(_{t\in I}\) of (B ⊗K B°)-modules isomorphic to B, and observe that for every \(t\in I\) , the element of M \(_{t}\) corresponding to 1 ∈ B by such an isomorphism belongs to B'.)

\begin{enumerate}
\def\labelenumi{\alph{enumi})}
\setcounter{enumi}{1}
\tightlist
\item
  Conversely, let B be a K-algebra such that for every algebra A containing B and having the same unit element as B, the canonical homomorphism from \(B \otimes_{K} B'\) to A is bijective. Prove that the algebra B is central simple and of finite rank (take \(A = \text{End}_{K}(B)\) ; deduce from the assumption that B is central and quasi-simple, and use Exercise 5, e) of VIII, p.~225).
\end{enumerate}

\begin{enumerate}
\def\labelenumi{\arabic{enumi})}
\setcounter{enumi}{3}
\item
  Let \(\mathrm{K}\) be a commutative field, \(\sigma\) an automorphism of \(\mathrm{K}\) distinct from the identity, and \(\mathrm{D}\) the field \(\mathrm{K}((\mathrm{X}))_{\sigma}\) (VIII, p.~19, Exercise 21, e)). Let \(\tau\) be an automorphism of \(\mathrm{K}\) that commutes with \(\sigma\) . Prove that there exists a unique automorphism of \(\mathrm{D}\) extending \(\tau\) and fixing X. Deduce an example of a field automorphism that fixes the center and is not inner.
\item
  Let K be a commutative field, V a vector space of finite dimension n over K, and W a subspace of V distinct from \(\{0\}\) and V. Denote by a the subset of the algebra \(A = \operatorname{End}_{K}(V)\) consisting of the endomorphisms u such that \(\operatorname{Im} u \subset W \subset \operatorname{Ker} u\) and by B the subalgebra \(K + a\) of A. Prove that B is a maximal commutative subalgebra of A, is not semisimple, and has degree possibly \textgreater n over K.
\end{enumerate}

¶ 6) Let D be a field with center K, distinct from K, such that every element of D is algebraic over K.

\begin{enumerate}
\def\labelenumi{\alph{enumi})}
\item
  Prove that D contains a separable commutative extension of degree \(>1\) .
\item
  Let m be an integer; suppose that every element of D has degree \(\leqslant m\) over K. Prove that we have \([D:K]\leqslant m^{2}\) (first prove that D has finite degree using a) and Theorem 5 of VIII, p.~259).
\end{enumerate}

\begin{enumerate}
\def\labelenumi{\arabic{enumi})}
\setcounter{enumi}{6}
\tightlist
\item
  Let A be a commutative ring, M and N A-modules, and \(u: M \to N\) an A-linear mapping. For every maximal ideal m of A, denote by \(u(\mathfrak{m})\) the homomorphism \(M/mM \to N/mN\) induced by u.
\end{enumerate}

\begin{enumerate}
\def\labelenumi{\alph{enumi})}
\item
  Suppose that N is finitely generated and that \(u(\mathfrak{m})\) is surjective for every maximal ideal m of A. Prove that u is surjective (apply Lemma 1 of VIII, p.~81 to the cokernel of u).
\item
  Suppose that M and N are finitely generated, that N is projective, and that \(u(\mathfrak{m})\) is bijective for every maximal ideal m of A. Prove that u is bijective (use the same method).
\item
  Suppose that M and N are projective and finitely generated and that \(u(\mathfrak{m})\) is injective for every maximal ideal m of A. Prove that u induces an isomorphism from M to a direct factor submodule of N (apply a) to the homomorphism \(^{t}u\) .
\item
  Let B be an A-algebra; suppose that the A-module B is faithful, projective, and finitely generated. Prove that A is a direct factor of the A-module B. Prove that if the B-module homomorphism \(u_{(B)} : M_{(B)} \to N_{(B)}\) is bijective, then so is u. Prove that if the B-module \(M_{(B)}\) is projective, then so is the A-module M.
\end{enumerate}

\begin{enumerate}
\def\labelenumi{\arabic{enumi})}
\setcounter{enumi}{7}
\tightlist
\item
  Let A be a commutative ring and S an A-algebra; suppose that S is a projective, faithful, and finitely generated A-module. Denote the A-algebra \(S \otimes_{A} S^{o}\) by E.
\end{enumerate}

\begin{enumerate}
\def\labelenumi{\alph{enumi})}
\tightlist
\item
  Prove that the following conditions are equivalent:
\end{enumerate}

\begin{enumerate}
\def\labelenumi{(\roman{enumi})}
\item
  The (E, A)-bimodule S is invertible.
\item
  The A-algebra S is central, and the E-module S is projective.
\item
  The canonical homomorphism \(E \longrightarrow \operatorname{End}_{\mathrm{A}}(\mathrm{S})\) is bijective.
\item
  For every maximal ideal m of A, the A/m-algebra S/mS is central simple. (Deduce the equivalence of (i), (ii), and (iii) from Theorem 1 of VIII, p.~101 and Corollary 2 of VIII, p.~100 and that of (iii) and (iv) from Theorem 1 of VIII, p.~252 and Exercise 7).
\end{enumerate}

When these conditions are satisfied, we say that \(S\) is an Azumaya algebra over \(A\) .

\begin{enumerate}
\def\labelenumi{\alph{enumi})}
\setcounter{enumi}{1}
\item
  If S and T are Azumaya algebras over A, then so is the A-algebra \(S \otimes_{A} T\) .
\item
  Let B be a commutative A-algebra; if S is an Azumaya algebra over A, then the B-algebra \(S_{(B)}\) is an Azumaya algebra. If the B-algebra \(S_{(B)}\) is an Azumaya algebra and B is a faithful and projective ( \(^{*}\) or faithfully flat \(_{*}\) ) A-algebra, then S is an Azumaya A-algebra (apply Exercise 7, d)).
\end{enumerate}

\begin{enumerate}
\def\labelenumi{\arabic{enumi})}
\setcounter{enumi}{8}
\tightlist
\item
  Let \(\mathrm{A}\) be a commutative ring and \(\mathrm{S}\) an Azumaya A-algebra (Exercise 8).
\end{enumerate}

\begin{enumerate}
\def\labelenumi{\alph{enumi})}
\item
  Prove that the center of S (which we identify with A) is a direct factor of the A-module S (use Exercise 7, d)).
\item
  Let M be an (S,S)-bimodule; denote by \(\mathcal{Z}(M)\) the A-submodule of M consisting of the elements m such that sm = ms for every \(s \in S\) . Deduce from Exercise 8, a) that the canonical homomorphism \(S \otimes_{A} \mathcal{Z}(M) \to M\) is bijective.
\item
  The mapping \(N \mapsto \mathcal{Z}(N)\) is an increasing bijection from the set of \((S, S)\) -sub-bimodules of M to the set of A-submodules of \(\mathcal{Z}(M)\) ; the inverse bijection sends an A-submodule N of \(\mathcal{Z}(M)\) to the \((S, S)\) -sub-bimodule SN of M. In particular, the mapping \(s \mapsto s \cap A\) is an increasing bijection from the set of two-sided ideals of S to the set of ideals of A.
\item
  Let \(\rho : \mathrm{S} \to \mathrm{T}\) be an A-algebra homomorphism, and let \(\mathrm{S}'\) be the commutant of \(\rho(\mathrm{S})\) in T. Prove that the canonical homomorphism \(\mathrm{S} \otimes_{\mathrm{A}} \mathrm{S}' \to \mathrm{T}\) is bijective (apply \(a\) ) to the (S,S)-bimodule T).
\item
  Suppose that every projective A-module L such that \(\dim_{\mathrm{A/m}}(\mathrm{L}/\mathfrak{m}\mathrm{L}) = 1\) for every maximal ideal m of A is free of rank 1 (this is, for example, the case when the ring A is local or a principal ideal domain). Prove that every endomorphism f of the A-algebra S is an inner automorphism (apply a) to the (S,S)-bimodule \(S^{f}\) , using Exercise 8, a) to show that the A-module \(\mathcal{Z}(S^{f})\) is projective).
\end{enumerate}

\begin{enumerate}
\def\labelenumi{\arabic{enumi})}
\setcounter{enumi}{9}
\tightlist
\item
  Let A be a commutative ring and P a faithful, projective, and finitely generated A-module.
\end{enumerate}

\begin{enumerate}
\def\labelenumi{\alph{enumi})}
\item
  Prove that the A-algebra \(\mathrm{S} = \operatorname{End}_{\mathrm{A}}(\mathrm{P})\) is an Azumaya algebra.
\item
  Prove that the left S-module P is projective.
\item
  Let \(P'\) be a second faithful, projective, and finitely generated A-module. Prove that the A-module \(P \otimes_{A} P'\) is faithful, projective, and finitely generated and that the algebra \(\text{End}_{\text{A}}(\text{P} \otimes_{\text{A}} \text{P}')\) can be canonically identified with the algebra \(\text{End}_{\text{A}}(\text{P}) \otimes_{\text{A}} \text{End}_{\text{A}}(\text{P}')\) .
\end{enumerate}

¶ 11) Let A be a commutative integral domain, K its field of fractions, and S an Azumaya A-algebra (Exercise 8).

\begin{enumerate}
\def\labelenumi{\alph{enumi})}
\item
  Prove that the K-algebra \(S_{(K)}\) is central simple. Identify S with a subalgebra of \(S_{(K)}\) .
\item
  Assume from now on that A is the only A-subalgebra of K that is a finitely generated A-module, \(^{*}\) that is, that the ring A is integrally closed, cf.~Comm. Alg., V, §1, No.~1, 2*. Prove that S is maximal (for the inclusion) among the A-subalgebras of \(S_{(K)}\) that are finitely generated as A-modules (apply Exercise 9, d)).
\item
  Suppose, moreover, that A is Noetherian and that the K-algebra \(S_{(K)}\) is isomorphic to \(\operatorname{End}_{\mathrm{K}}(\mathrm{V})\) , where V is a vector space over K. Prove that there exists a finitely generated A-module M without torsion such that the A-algebra S is isomorphic to \(\operatorname{End}_{\mathrm{A}}(\mathrm{M})\) (let x be a nonzero vector of V; take M to be the set of elements \(s(x)\) as s runs through S).
\end{enumerate}

\(^{*}d\) Prove that in the previous construction, we may assume M reflexive (replace M with its bidual). If, moreover, the ring A is regular (AC, X, §4, n° 2, p.~55, définition 1), then the A-module M is projective (apply AC, X, §4, p.~164, exercice 12).*

*¶ 12) Let K be a commutative field, complete for a discrete valuation v (Comm. Alg., VI) that we assume normalized (that is, \(v(\mathrm{K}^{*}) = \mathbf{Z}\) ). Denote the ring of v by A and its residue field by \(\kappa\) . Recall (loc. cit.) that if L is a finite extension of K, then the valuation v admits a unique extension \(v_{L}\) to L; we say that the extension is unramified if \(v_{L}\) has integer values.

\begin{enumerate}
\def\labelenumi{\alph{enumi})}
\item
  Let D be a field with center K; let \(Nrd : D \to K\) be the reduced norm on D (VIII, p.~340, Definition 2). Prove that the mapping \(w_{0} = v \circ Nrd\) is a valuation on D (to prove the inequality \(w_{0}(x + y) \geqslant \inf(w_{0}(x), w_{0}(y))\) , reduce to the case y = 1, and observe that the restriction of \(w_{0}\) to \(K(x)\) is a multiple of \(v_{\mathrm{K}(x)}\) ). Denote by w the normalized valuation on D equivalent to \(w_{0}\) .
\item
  An element x of L is integral over A if and only if we have \(w(x) \geqslant 0\) ; the set of these elements forms a subring B of D.
\item
  Assume from now on that the field \(\kappa\) is perfect. Prove that if D is distinct from K, then it contains an unramified extension of K distinct from K. (In the opposite case, observe that we have \([\kappa_{\mathrm{L}}:\kappa] = 1\) for every extension L of K contained in D. Let \(b\in \mathrm{B}\) , and let \(\pi\) be a uniformizer of D. Apply the previous remark to the field \(\mathrm{K}(b)\) , and prove that \(b\) is in the closure of \(\mathrm{A}(\pi)\) . Since the subspace \(\mathrm{K}(\pi)\) is closed in D (EVT, I, §2, n° 3, p.~14, corollaire 1), we obtain \(\mathrm{B}\subset \mathrm{K}(\pi)\) and finally \(\mathrm{D} = \mathrm{K}(\pi)\) .) d) Prove that there exists a maximal commutative subfield of D that is unramified over K. (Use induction on the reduced degree of D by considering a subfield L of D that is unramified over K and applying the induction hypothesis to the commutant of L in D.)*
\end{enumerate}

\begin{enumerate}
\def\labelenumi{\arabic{enumi})}
\setcounter{enumi}{12}
\tightlist
\item
  Let D be a field, V a left vector space over D, A the ring \(\operatorname{End}_{\mathrm{D}}(\mathrm{V})\) , and f an automorphism of A. Prove that there exist an automorphism \(\sigma\) of D and a bijective mapping \(u: V \to V\) that is semilinear with respect to \(\sigma\) (II, §1, No.~13, p.~223) such that we have \(f(a) = u \circ a \circ u^{-1}\) for every \(a \in A\) . (Let \(V^{f}\) be the A-module with additive group V and external law defined by \((a, v) \mapsto f(a)(v)\) ; deduce from
\end{enumerate}

Proposition 4 of VIII, p.~50 an A-isomorphism \(u: V \to V^{f}\) . Then, observe that \(uDu^{-1}\) is equal to the bicommutant of the D-module V, that is, to D.)

\begin{enumerate}
\def\labelenumi{\arabic{enumi})}
\setcounter{enumi}{13}
\tightlist
\item
  Let D be a field, Z its center, V a left vector space over D, and \(\Omega\) the endomorphism ring of the abelian group V. Endow \(\Omega\) with its natural (D,D)-bimodule structure. Let \(\Gamma\) be the automorphism group of D and \(\Delta\) the subgroup of inner automorphisms (isomorphic to \(D^{*}/Z^{*}\) ).
\end{enumerate}

Let S be a subset of \(\Omega\) consisting of semilinear mappings (relative to an element of \(\Gamma\) ) and F the left D-linear subspace of \(\Omega\) generated by S; it coincides with the right D-linear subspace generated by S.

\begin{enumerate}
\def\labelenumi{\alph{enumi})}
\item
  For \(\theta \in \Gamma / \Delta\) , denote by \(S_{\theta}\) the set of semilinear mappings from V to itself with respect to an automorphism of D belonging to \(\theta\) and by \(F_{\theta}\) the D-linear space of F generated by \(S_{\theta}\) . Prove that F is the direct sum of the subspaces \(F_{\theta}\) for \(\theta \in \Gamma / \Delta\) . b) Prove that if \(S_{\theta}\) is a free subset over Z, then it is also (left and right) free over D (reason as in the proof of Theorem 1 of V, §6, No.~1, p.~27).
\item
  Let E be a (D,D)-sub-bimodule of F. Prove that E is generated (over D) by semilinear mappings belonging to the \(F_{\theta}\) . (Let B be a basis of F contained in S. Prove that E is generated by the vectors \(u \in E\) that have the following property: if \(u = \sum_{i}^{n} \lambda_{i} b_{i}\) with \(\lambda_{i} \neq 0\) and \(b_{i} \in B\) for every i, then the intersection of E and the subspace of F generated by the \(b_{i}\) is reduced to Du. Then deduce from a) that u is semilinear.)
\end{enumerate}

¶ 15) Keep the notation of Exercise 14; denote the ring \(\mathrm{End}_{\mathrm{D}}(\mathrm{V})\) by A. Let G be a group of automorphisms of A and \(\mathrm{G}_0\) the subgroup of elements of G that are inner automorphisms of A. Denote by U(G) the (left and right) D-linear subspace of \(\Omega\) generated by the semilinear bijections \(s\) of V such that the automorphism \(a \mapsto sas^{-1}\) belongs to G, and denote by \(\mathrm{U}_0(\mathrm{G})\) the Z-linear subspace of A generated by the automorphisms \(u\) of V such that the automorphism \(a \mapsto uau^{-1}\) belongs to \(\mathrm{G}_0\) . These are Z-subalgebras of \(\Omega\) .

\begin{enumerate}
\def\labelenumi{\alph{enumi})}
\item
  Let \((s_{\alpha})\) be a basis of \(U_{0}(G)\) over Z consisting of elements of \(S_{0}(G)\) , and let \((g_{\beta})_{\beta\in G/G_{0}}\) be a system of representatives of the classes (mod \(G_{0}\) ) in G. For each \(\beta\in G/G_{0}\) , choose an element \(u_{\beta}\) of \(S(G)\) corresponding to \(g_{\beta}\) . Prove that the elements \(s_{\alpha}u_{\beta}\) form a (left and right) basis of \(U(G)\) over D (use Exercise 14, a)). The degree \([U(G):D]\) is finite if and only if \([U_{0}(G):Z]\) and \((G:G_{0})\) are finite; if this is the case, then we have \([U(G):D]=[U_{0}(G):Z](G:G_{0})\) .
\item
  Prove that \(\mathrm{U}_{0}(\mathrm{G})\) is equal to \(\mathrm{U}(\mathrm{G})\cap\mathrm{A}\) and that \(\mathrm{U}(\mathrm{G}_{0})=\mathrm{D}u_{0}(\mathrm{G})\) is canonically isomorphic to \(\mathrm{U}_{0}(\mathrm{G})\otimes_{\mathrm{Z}}\mathrm{D}\) (use Exercise 14 as well as Exercise 14 of VIII, p.~227). c) An element of \(\mathrm{U}(\mathrm{G})\) is a semilinear mapping from V to itself if and only if it is of the form \(\lambda su_{\beta}\) with \(\lambda\in D\) , \(s\in\mathrm{U}_{0}(\mathrm{G})\) , and \(\beta\in\mathrm{G}/\mathrm{G}_{0}\) (use a), b), and Exercise 14, a)).
\item
  Denote the commutant of U(G) in \(\Omega\) by \(A^{G}\) ; it is the subring of A consisting of the elements invariant under G. Set \(Z^{G} = A^{G} \cap Z\) ; it is the subfield of elements of Z invariant under G. Prove that Z and \(A^{G}\) are linearly disjoint over \(Z^{G}\) (verify that a basis of \(A^{G}\) over \(Z^{G}\) is a free subset of A over Z).
\item
  Prove that if \(U_{0}(G)\) is a simple ring and one of the Z-algebras \(U_{0}(G)\) and D has finite degree, then the ring \(U(G)\) is quasi-simple (deduce from c) that \(U(G_{0})\) is a simple ring, and apply Exercise 14, c) to a two-sided ideal of \(U(G)\) using c)).
\item
  Suppose that V and U(G) are finite-dimensional over D and that the ring \(U_{0}(G)\) is simple. Prove that the ring A{[}G{]} is simple and that we have \([A : A^{G}]_{s} = [U_{0}(G) : Z](G : G_{0})\) (apply Proposition 12 of VIII, p.~205). Prove that the commutant of \(A^{G}\) in \(\Omega\) (resp. in A) is U(G) (resp. \(U_{0}(G)\) ) (use Proposition 5 of VIII, p.~139).
\end{enumerate}

\(g\) ) Deduce that the group of automorphisms of A that fix \(\mathrm{A}^{\mathrm{G}}\) consists of the automorphisms \(a\mapsto vav^{-1}\) , where \(v\) runs through the set of semilinear bijections in \(\mathrm{U(G)}\) . Prove that it is generated by G and the inner automorphisms associated with the invertible elements of \(\mathrm{U_0(G)}\) (use \(c\) ).

\begin{enumerate}
\def\labelenumi{\alph{enumi})}
\setcounter{enumi}{7}
\tightlist
\item
  If \(G = G_{0}\) , then the ring \(A^{G}\) contains Z; prove the converse when D is finite-dimensional over Z (use Theorem 4 of VIII, p.~257).
\end{enumerate}

¶ 16) Keep the notation of Exercise 15, and suppose, moreover, that V is finite-dimensional over D. A subring B of A is called weakly Galois (resp. Galois) in A if its commutant B' in Ω is generated as a (left or right) D-vector space by semilinear mappings (resp. bijections) from V to itself. Denote by DB the subring of Ω generated by D and B.

\begin{enumerate}
\def\labelenumi{\alph{enumi})}
\item
  Prove that if B is simple and weakly Galois, then V is a semisimple module of finite length over DB. (Let W be an isotypical component of the B-module V and S a simple DB-submodule of W. Let \(\Sigma\) be the set of semilinear mappings v belonging to \(B'\) . Observe that \(v(\mathrm{S})\) is an DB-module for \(v \in \Sigma\) and deduce from Corollary 3 of VIII, p.~56 that W is contained in the sum of the \(v(\mathrm{S})\) for \(v \in \Sigma\) .)
\item
  Suppose that the ring B is simple and weakly Galois and that its commutant \(B_{0}^{\prime}\) in A is simple. Prove that V is an isotypical DB-module (observe that \(B_{0}^{\prime}\) is the commutant of DB in \(\Omega\) , and use Proposition 5 of VIII, p.~85). Deduce that B is Galois. (Let v be a nonzero element of \(\Sigma\) and \(V = \bigoplus_{i=1}^{n} V_{i}\) a decomposition of V as a direct sum of isomorphic simple DB-modules such that \(W_{1} = v(V_{1})\) is nonzero. Observe that there exist simple DB-modules \(W_{i}\) (for \(2 \leqslant i \leqslant n\) ) isomorphic to \(W_{1}\) such that \(V = \bigoplus_{i=1}^{n} W_{i}\) , and deduce a semilinear bijection \(w \in B^{\prime}\) that coincides with v on \(V_{1}\) and for which \(vw^{-1}\) belongs to \(B_{0}^{\prime}\) . Conclude by observing that every element of a simple ring is a sum of invertible elements.)
\end{enumerate}

¶ 17) Keep the notation of Exercise 15, and suppose that V is finite-dimensional over D. A group G of automorphisms of A is called Galois if it satisfies the following conditions:

\begin{enumerate}
\def\labelenumi{(\roman{enumi})}
\item
  \((\mathrm{G}:\mathrm{G}_{0})\) and \([\mathrm{U}_{0}(\mathrm{G}):\mathrm{Z}]\) are finite.
\item
  The ring \(U_{0}(G)\) is simple.
\item
  For every invertible element s of \(\mathrm{U}_{0}(\mathrm{G})\) , the automorphism \(a \mapsto sas^{-1}\) belongs to \(G_{0}\) .
\end{enumerate}

Let \(G\) be such a group.

\begin{enumerate}
\def\labelenumi{\alph{enumi})}
\item
  Prove that the mapping \(H \mapsto A^{H}\) is a decreasing bijection from the set of Galois subgroups of G to the set of simple subrings of A that contain \(A^{G}\) and whose commutant \(B_{0}^{\prime}\) in A is simple. (Deduce from Exercise 15, f) and g) that this mapping is well defined and injective. Let B be a simple subring of A satisfying the conditions above, and let H be the group of automorphisms of A fixing B. Use Exercises 14, c) and 16 to prove that B is Galois and then Proposition 12 of VIII, p.~205 to prove that we have \(B = A^{H}\) and that H is Galois.)
\item
  Let \(B_{1}\) and \(B_{2}\) be simple subrings of A containing \(A^{G}\) whose commutants in A are simple, and let \(\varphi\) be an isomorphism from \(B_{1}\) to \(B_{2}\) fixing \(A^{G}\) . Prove that there exists an element of G that coincides with \(\varphi\) on \(B_{1}\) .
\end{enumerate}

(First observe that we have \(\operatorname{long}_{\mathrm{B}_{1}}(\mathrm{V}) = \operatorname{long}_{\mathrm{B}_{2}}(\mathrm{V})\) , using the equality \(h(\mathrm{B}_{1}, \mathrm{A}^{\mathrm{G}}) = h(\mathrm{B}_{2}, \mathrm{A}^{\mathrm{G}})\) . Let W be the set of elements f of \(\Omega\) satisfying \(f(bx) = \varphi(b)f(x)\) for \(b \in B_{1}\) , \(x \in V\) . Prove that W is not reduced to 0, and then apply Exercise 14, c) to it to show that it contains a nonzero semilinear mapping w. Show that if S is a simple \(DB_{1}\) -submodule of V that is not contained in Ker w, then w induces a bijection from S to its image and we have \(\operatorname{long}_{\mathrm{B}_{1}}(\mathrm{S}) = \operatorname{long}_{\mathrm{B}_{2}}(w(\mathrm{S}))\) (use Exercise 16, b)). Reasoning as in loc. cit., conclude that W contains a semilinear bijection that coincides with w on S.)

\begin{enumerate}
\def\labelenumi{\alph{enumi})}
\setcounter{enumi}{2}
\tightlist
\item
  Specialize the results of a) and b) to the following cases:
\end{enumerate}

\begin{enumerate}
\def\labelenumi{(\roman{enumi})}
\item
  \(G = G_{0}\) (cf.~VIII, p.~259, Theorem 5).
\item
  \(\mathrm{G}_0 = \{1\}\) . (Prove that every finite group of automorphisms of A is Galois and that every simple subring of A containing \(\mathrm{A}^{\mathrm{G}}\) admits Z as commutant in A.)
\item
  V has dimension 1, so that A is a field isomorphic to \(D^{o}\) . (The ring \(A^{G}\) is a field, and every subfield of A containing \(A^{G}\) has a field as commutant. Study the specific case when D is commutative (cf.~V, §10, No.~7).)
\end{enumerate}

\begin{enumerate}
\def\labelenumi{\arabic{enumi})}
\setcounter{enumi}{17}
\tightlist
\item
  Let \(\mathrm{K}\) be a commutative field, \(\mathbf{Z}\) a cyclic extension of \(\mathrm{K}\) (V, §11, No.~6, p.~85) of degree \(n \geqslant 3\) , and \(\sigma\) a generator of the Galois group of \(\mathbf{Z}\) over \(\mathbf{K}\) . Denote by A the ring \(\mathbf{M}_2(\mathbf{Z})\) and by B the subring of A consisting of the matrices \(\left( \begin{array}{cc}z & 0\\ 0 & \sigma (z) \end{array} \right)\) for \(z\in \mathbf{Z}\) . a) Keep the notation of Exercise 17. Prove that there exists a Galois group \(\mathbf{G}\) of automorphisms of A such that \(\mathrm{A}^{\mathrm{G}} = \mathrm{K}\) , \((\mathrm{G}:\mathrm{G}_0) = n\) , and \(\mathrm{U}_0(\mathrm{G}) = \mathrm{A}\) .
\end{enumerate}

\begin{enumerate}
\def\labelenumi{\alph{enumi})}
\setcounter{enumi}{1}
\item
  Prove that B is a simple subring of A but that its commutant \(B_{0}^{\prime}\) in A is isomorphic to \(Z \times Z\) .
\item
  Let H be the group of automorphisms of A fixing B. Prove that \(A^{H}\) is equal to \(B_{0}^{\prime}\) . Construct an isomorphism from B to Z that fixes K but does not extend to an automorphism of A (cf.~Exercise 17, b)).
\item
  Prove that the subring B of A is weakly Galois but not Galois (Exercise 16; describe the commutant of B in \(\Omega\) ).
\end{enumerate}

¶ 19) Let D be a field, Z its center, and V a vector space over D. Identify D with the subfield of homotheties of V. Let A be a dense subring of \(\mathrm{End}_{\mathrm{D}}(\mathrm{V})\) (VIII, p.~129, Exercise 3) containing Z. Let E be a subfield of D containing Z and F its commutant in D. Prove that the Z-algebra \(\mathrm{A} \otimes_{\mathrm{Z}} \mathrm{E}\) is isomorphic to a dense subalgebra of \(\mathrm{End}_{\mathrm{F}}(\mathrm{V})\) (use Exercise 14 of VIII, p.~227, and consider the commutant of AE in \(\mathrm{End}_{\mathrm{Z}}(\mathrm{V})\) ). This is the case when E is a maximal subfield of D.

